% from ega2-1-fr.tex
\oldpage[II]{166}
à l'injection $S_f^{\leq}\to S_f$ (\hyperref[II.8.3.2-fr]{8.3.2}). Enfin, la
dernière assertion est conséquence de ce que pour $f\in S_1$, $S_f$ est
isomorphe à $S_{(f)}[T,T^{-1}]$ (\hyperref[II.2.2.1-fr]{2.2.1}).

\begin{remark}[8.3.7]
\label{II.8.3.7-fr}
Lorsque $Y$ est affine, les éléments de l'espace sous-jacent à $E$ sont les
idéaux premiers $\mathfrak{p}$ (non nécessairement gradués) de $S$ ne
contenant pas $S_+$, en vertu de la définition de l'immersion $\varepsilon$
(\hyperref[II.8.3.2-fr]{8.3.2}). Pour un tel idéal $\mathfrak{p}$, les
$\mathfrak{p}\cap S_n$ vérifient de façon évidente les conditions de
(\hyperref[II.2.1.9-fr]{2.1.9}), donc il existe un seul idéal premier
\emph{gradué} $\mathfrak{q}$ de $S$ tel que
$\mathfrak{q}\cap S_n=\mathfrak{p}\cap S_n$ pour tout $n$~; l'application
$\pi:E\to X$ des espaces sous-jacents s'interprète alors grâce à la relation
\begin{equation}
\label{II.8.3.7.1-fr}
  \pi(\mathfrak{p})=\mathfrak{q}.
\tag{8.3.7.1}
\end{equation}
En effet, pour vérifier cette relation, il suffit de considérer un $f$
homogène dans $S_+$ tel que $\mathfrak{p}\in D(f)$, et d'observer que
$\mathfrak{q}_{(f)}$ est l'image réciproque de $\mathfrak{p}_f$ par
l'injection $S_{(f)}\to S_f$.
\end{remark}

\begin{corollary}[8.3.8]
\label{II.8.3.8-fr}
Si $\mathcal{S}$ est engendrée par $\mathcal{S}_1$, les morphismes $p$ et
$\pi$ sont de type fini~; pour tout $x\in X$, la fibre $p^{-1}(x)$ est
isomorphe à $\operatorname{Spec}(k(x)[T])$ et la fibre $\pi^{-1}(x)$ est
isomorphe à $\operatorname{Spec}(k(x)[T,T^{-1}])$.
\end{corollary}

Cela résulte aussitôt de (\hyperref[II.8.3.5-fr]{8.3.5}) et
(\hyperref[II.8.3.6-fr]{8.3.6}) en observant que, lorsque $Y$ est affine et
$S$ engendrée par $S_1$, les $X_f$, pour $f\in S_1$, forment un recouvrement
de $X$ (\hyperref[II.2.3.14-fr]{2.3.14}).

\begin{remark}[8.3.9]
\label{II.8.3.9-fr}
Le cône épointé affine correspondant à la $\mathcal{O}_Y$-Algèbre graduée
$\mathcal{O}_Y[T]$ ($T$ indéterminée) s'identifie à
$G_m=\operatorname{Spec}(\mathcal{O}_Y[T,T^{-1}])$, puisqu'il n'est autre que
$C_T$, comme on l'a vu dans (\hyperref[II.8.3.2-fr]{8.3.2}) (voir
(\hyperref[II.8.4.4-fr]{8.4.4}) pour un résultat plus général). Ce préschéma
est canoniquement muni d'une structure de «~\emph{$Y$-schéma en groupes
commutatif}~». Il s'agit là d'une notion qui sera développée en détail plus
tard, et que l'on peut ici résumer rapidement comme suit. Un $Y$-schéma en
groupes est un $Y$-schéma $G$, muni de deux $Y$-morphismes
$p:G\times_YG\to G$ et $s:G\to G$ qui satisfont aux conditions formellement
analogues aux axiomes de la loi de composition et de l'application de
symétrie dans un groupe~: le diagramme
\[
  \xymatrix{
    G\times G\times G \ar[r]^{p\times1}\ar[d]_{1\times p}
      & G\times G \ar[d]^{p} \\
    G\times G \ar[r]_{p} & G
  }
\]
doit être commutatif («~associativité~») et l'on doit avoir une condition qui
correspond dans les groupes au fait que les applications
\[
  (x,y)\mapsto(x,x^{-1},y)\mapsto(x,x^{-1}y)\mapsto x(x^{-1}y)
\]
et
\[
  (x,y)\mapsto(x,x^{-1},y)\mapsto(x,yx^{-1})\mapsto(yx^{-1})x
\]
doivent toutes deux se réduire à $(x,y)\mapsto y$~; la suite de morphismes
correspondant par exemple à la première application composée est
\[
  G\times G\xrightarrow{(1,s)\times1}G\times G\times G
  \xrightarrow{1\times p}G\times G\xrightarrow{p}G
\]
et le lecteur écrira de même la seconde suite.

